\mxunit{premiere-c04}{Dérivation : du taux à la tangente}\label{C04}
\begin{keypoint}[Le parcours du chapitre]
\textbf{Objectifs.} Interpréter un taux ; approcher un nombre dérivé ; déterminer une tangente ; comprendre les échecs de dérivabilité.

\textbf{Prérequis.} Images, pente d’une droite, calcul littéral, racine carrée et valeur absolue.

\textbf{Repère annuel.} 3–12 novembre : 8 h\end{keypoint}
\section*{Diagnostic : \mxcount{diag} questions, 10 minutes}\label{C04-D}
Répondre sans consulter le cours. Indiquer les questions pour lesquelles une aide est nécessaire. Ce diagnostic ne donne pas lieu à une note.
\begin{enumerate}
\item Pente de la droite passant par $\pt{1;2}$ et $\pt{3;8}$ ?
\item Développer $(2+h)^2-4$.
\item Simplifier $\dfrac{4h+h^2}{h}$ si $h\ne0$.
\item Calculer $\sqrt0$ et $\sqrt{\num{0,01}}$.
\item Calculer $|-3|$ et $|2|$.
\item Une droite de pente 2 passe par $\pt{1;5}$. Équation ?
\mxcounted{diag}
\end{enumerate}
\paragraph{Après le diagnostic.} Reprendre les notions signalées avec le professeur ; refaire une question analogue, puis expliquer la règle utilisée.
\paragraph{Oral et logique.} Expliquer à un camarade pourquoi on simplifie le taux avant de faire tendre h vers 0. Distinguer vitesse moyenne et instantanée.
\mxnewpage
\section*{Chercher avant de formaliser}\label{C04-RECH}
Pour $f(x)=x^2$, faire varier h positif et négatif autour de 0 en gardant le point A d’abscisse 1. Dessiner trois sécantes ; conjecturer une droite limite puis la déterminer sans lecture graphique.
\paragraph{Organisation.} Recherche individuelle 5 min, confrontation à deux 5 min, mise en commun 5 min. Garder une trace d’un essai, même non abouti. Une conjecture doit ensuite être justifiée.
\mxnewpage
\section{Taux de variation et sécantes}\label{C04-U1}
\begin{keypoint}[Définition et résultats]
Pour $a\ne b$ dans le domaine de f, le taux de variation entre a et b vaut $\dfrac{f(b)-f(a)}{b-a}$. C’est la pente de la sécante passant par les points correspondants. En contexte, l’unité du taux est l’unité de f divisée par celle de la variable.
\end{keypoint}
\paragraph{Statut et justification.} Définition ; formule de la pente d’une droite déjà étudiée. Avec $b=a+h$, $h\ne0$, le taux devient $\dfrac{[f(a+h)-f(a)]}{h}$.
\begin{example}Pour $f(x)=x^2$, entre 1 et 3, le taux vaut $\dfrac{9-1}{3-1}=4$. Entre 1 et $1+h$, il vaut $2+h$. Les pentes pour $h=\num{0,1}$ et $h=-\num{0,1}$ sont \num{2,1} et \num{1,9}.\end{example}\label{C04-EX1}
\begin{pitfall}Le taux n’est pas $\dfrac{f(b)}{f(a)}$ ; h représente un écart, positif ou négatif selon le domaine.\end{pitfall}
\begin{application}Calculer le taux de $f(x)=x^2+3x$ entre 2 et $2+h$.\end{application}\label{C04-AP1}
\section{Nombre dérivé}\label{C04-U2}
\begin{keypoint}[Définition et résultats]
f est dérivable en a lorsque le taux $\dfrac{[f(a+h)-f(a)]}{h}$ tend vers un réel quand h tend vers 0, dans le domaine. Ce réel est $f\prime(a)$. On emploie ici une approche intuitive de la limite. En cinématique, si f donne une position en mètres et le temps en secondes, le nombre dérivé représente une vitesse instantanée en $\mathrm{m\,s^{-1}}$.
\end{keypoint}
\paragraph{Statut et justification.} Pour $f(x)=x^2$, le taux en a vaut $2a+h$, qui se rapproche de $2a$ ; ainsi $f\prime(a)=2a$. Ce calcul sera repris comme fonction dérivée au chapitre 6.
\begin{example}Pour $s(t)=t^2+3t$, le taux autour de 2 vaut $7+h$, donc $s\prime(2)=7\ \mathrm{m\,s^{-1}}$. La vitesse moyenne sur $\intff{2}{3}$ vaut $8\ \mathrm{m\,s^{-1}}$ ; elle n’est pas la vitesse à l’instant 2.\end{example}\label{C04-EX2}
\begin{pitfall}On ne remplace pas h par 0 dans un quotient non simplifié : on obtiendrait une division par 0.\end{pitfall}
\begin{application}Pour $f(x)=3x^2-2x$, calculer $f\prime(1)$ à partir du taux.\end{application}\label{C04-AP2}
\section{Tangente et approximation affine}\label{C04-U3}
\begin{keypoint}[Définition et résultats]
Si f est dérivable en a, la tangente au point $\pt{a;f(a)}$ a pour équation $y=f(a)+f\prime(a)(x-a)$. Pour h petit, on utilise $f(a+h)\approx f(a)+f\prime(a)h$. Cette approximation locale ne fournit pas à elle seule une borne de l’erreur.
\end{keypoint}
\paragraph{Statut et justification.} La tangente a pente $f\prime(a)$ et passe par $\pt{a;f(a)}$. Dans $y=f\prime(a)x+p$, substituer ce point donne $p=f(a)-af\prime(a)$, d’où l’équation.
\begin{example}Pour $f(x)=x^2$ en 2, tangente $y=4+4(x-2)=4x-4$. On approche $\num{2,01}^2$ par $4+4(\num{0,01})=\num{4,04}$ ; la valeur exacte est $\num{4,0401}$, écart $\num{0,0001}=h^2$.\end{example}\label{C04-EX3}
\mxinclude{figures/C04-U3}
\begin{pitfall}Une tangente peut recouper la courbe ailleurs. « Un seul point commun » n’est pas une définition.\end{pitfall}
\begin{application}On donne $f(3)=5$ et $f\prime(3)=-2$. Écrire la tangente et approcher $f(\num{3,04})$.\end{application}\label{C04-AP3}
\section{Quand la dérivabilité échoue}\label{C04-U4}
\begin{keypoint}[Définition et résultats]
La fonction racine carrée est définie sur $\intfo{0}{+\infty}$ mais n’est pas dérivable en 0. La fonction valeur absolue est définie sur $\R$ mais n’est pas dérivable en 0. Être définie ou avoir une courbe continue ne garantit donc pas la dérivabilité.
\end{keypoint}
\paragraph{Statut et justification.} Pour $h>0$, $\dfrac{\sqrt h-0}{h}=\dfrac{1}{\sqrt h}$, qui ne tend pas vers un réel fini lorsque h se rapproche de 0. Pour $|x|$, le taux $\dfrac{|h|}{h}$ vaut 1 à droite et $-1$ à gauche : pas de valeur commune.
\begin{example}Pour $|x|$, les sécantes issues de l’origine ont toutes pente 1 à droite et $-1$ à gauche. Le point anguleux ne possède pas une tangente de pente finie unique. Pour la racine, la direction limite est verticale ; il n’y a pas de nombre dérivé réel.\end{example}\label{C04-EX4}
\begin{pitfall}Une pente très grande n’est pas une preuve de non-dérivabilité ; il faut expliquer le comportement du taux.\end{pitfall}
\begin{application}Justifier que $g(x)=|x-2|$ n’est pas dérivable en 2. Étudier aussi la dérivabilité de $x\mapsto\sqrt{x}$ en 0 par son taux.\end{application}\label{C04-AP4}
\mxnewpage
\section*{Parcours d’exercices}\label{C04-PARC}
Commencer par les exercices 1 et 2. Choisir ensuite un exercice 3 ou 4 sans aide. Les exercices 5 à 7 sont une réserve de transfert, de recherche ou de programmation : le professeur en désigne un. L’exercice 8 se reprend à distance.
\begingroup
\mxexercisehead{C04-E1}{Deux sécantes}
La fonction $f$ est définie sur $\R$ par $f(x)=x^2-1$. On note $A$ le point de sa courbe d’abscisse 1.
\begin{enumerate}[label=\arabic*.]
\item Calculer le coefficient directeur de la sécante joignant $A$ au point de la courbe d’abscisse 2.
\item Reprendre le calcul avec le point de la courbe d’abscisse $\num{1,1}$.
\item Après l’étude du nombre dérivé, déterminer $f\prime(1)$ et comparer cette valeur aux deux taux calculés. Expliquer ce qui change lorsque le second point se rapproche de $A$.
\end{enumerate}
\mxexercisehead{C04-E2}{Une tangente}
On considère la fonction $f(x)=2x^2+1$, définie sur $\R$. L’objectif est de déterminer l’équation de sa tangente au point d’abscisse $-1$, en partant du taux de variation.
\begin{enumerate}[label=\arabic*.]
\item Pour $h\ne0$, écrire puis simplifier $\dfrac{f(-1+h)-f(-1)}{h}$.
\item En déduire $f\prime(-1)$. Calculer également $f(-1)$.
\item Écrire l’équation réduite de la tangente et vérifier qu’elle passe par le point de contact.
\end{enumerate}
\mxexercisehead{C04-E3}{Coût marginal}
Une entreprise utilise le modèle $C(q)=q^2+20q+100$, où $q\geq0$ désigne la quantité produite et $C(q)$ son coût en unités monétaires. Le modèle prolonge la quantité par une variable réelle afin d’étudier de petites variations de production.
\begin{enumerate}[label=\arabic*.]
\item Calculer $C(10)$ et $C(11)$, puis le taux de variation de $C$ entre 10 et 11. Interpréter ce taux dans le contexte.
\item À partir du taux entre 10 et $10+h$, déterminer $C\prime(10)$.
\item Comparer les deux résultats. Expliquer pourquoi le coût d’une unité supplémentaire n’est pas exactement égal au nombre dérivé dans ce modèle.
\end{enumerate}
\mxexercisehead{C04-E4}{Lire une tangente}
Cet exercice comporte deux situations indépendantes. Une tangente à une courbe est une droite : un autre point de cette droite n’appartient pas nécessairement à la courbe.
\begin{enumerate}[label=\arabic*.]
\item Sur la figure de la section consacrée à la tangente, lire les coordonnées de deux points de la droite verte. En déduire sa pente et vérifier la valeur annoncée dans la légende.
\item Dans une autre situation, la tangente à la courbe d’une fonction $f$ au point $A\pt{2;3}$ passe par $B\pt{4;2}$. Déterminer $f(2)$ et $f\prime(2)$. Peut-on en déduire $f(4)$ ?
\item Utiliser cette tangente pour proposer une valeur approchée de $f(\num{2,1})$. Faire apparaître le signe « environ égal » dans la conclusion.
\end{enumerate}
\mxexercisehead{C04-E5}{Contrôler une approximation}
On souhaite approcher $\num{0,98}^2$ sans effectuer d’abord la multiplication. On utilise la tangente à la courbe de la fonction carré au point d’abscisse 1.
\begin{enumerate}[label=\arabic*.]
\item Écrire l’équation de cette tangente et en déduire l’approximation recherchée.
\item Calculer ensuite la valeur exacte de $\num{0,98}^2$. On définit l’erreur par « valeur exacte moins valeur approchée » : calculer cette erreur et préciser son signe.
\item Expliquer le résultat à partir de l’identité $(1+h)^2=1+2h+h^2$.
\end{enumerate}
\mxexercisehead{C04-E6}{Tangente et intersection}
La fonction $f$ est définie sur $\R$ par $f(x)=x^3$. On admet que $f\prime(1)=3$. Un élève affirme qu’une tangente et une courbe ont toujours un seul point commun.
\begin{enumerate}[label=\arabic*.]
\item Déterminer l’équation de la tangente à la courbe de $f$ au point d’abscisse 1.
\item Vérifier en développant que $x^3-3x+2=(x-1)^2(x+2)$.
\item Déterminer tous les points d’intersection de la courbe et de cette tangente. L’affirmation de l’élève est-elle correcte ?
\end{enumerate}
\mxexercisehead{C04-E7}{Sécantes numériques}
Pour étudier la fonction carré au voisinage de 1, on définit, pour $h\ne0$,\[T(h)=\dfrac{(1+h)^2-1}{h}.\]On utilisera les six valeurs $\num{0,1}$, $\num{0,01}$, $\num{0,001}$ et leurs opposés.
\begin{enumerate}[label=\arabic*.]
\item Écrire une fonction Python qui calcule $T(h)$, puis produire un tableau de six résultats.
\item Conjecturer la valeur du nombre dérivé de la fonction carré en 1.
\item Simplifier algébriquement $T(h)$ pour justifier la conjecture. Expliquer pourquoi l’essai de valeurs numériques ne suffit pas, à lui seul, à démontrer le résultat.
\end{enumerate}
\mxexercisehead{C04-E8}{Reprise différée}
La fonction inverse est définie par $f(x)=\dfrac1x$ pour $x\ne0$. On l’étudie au voisinage de 1.
\begin{enumerate}[label=\arabic*.]
\item Préciser les valeurs interdites de $h$ dans le taux $\dfrac{f(1+h)-f(1)}h$, puis le simplifier.
\item En déduire $f\prime(1)$ et l’équation de la tangente en 1.
\item Utiliser cette tangente pour approcher $\dfrac1{\num{1,02}}$. Conserver la distinction entre valeur exacte et valeur approchée.
\end{enumerate}
\paragraph{Travail personnel.} Effectuer uniquement la tâche désignée à l’issue de la séance, dans la progression datée : 10–20 min maximum. Les exercices dépendent des notions déjà étudiées ; certaines questions sont donc différées. Reprendre l’exercice 8 une semaine après le chapitre (10 min). Si le temps prévu est dépassé, noter la difficulté et apporter les essais ; ne pas accumuler toute la réserve.
\endgroup
\mxneed{12}\section*{À faire dans le manuel CQFD}
Ces exercices se traitent dans votre manuel. La progression indique la séance ; ils remplacent le travail personnel prévu ce jour-là (10 à 20 minutes au total).
\paragraph{CQFD, exercice 17 p. 127, questions toutes.}\label{C04-M1} Lire deux points de chaque tangente ; écrire le quotient des accroissements plutôt qu’estimer son inclinaison.
\paragraph{CQFD, exercice 19 p. 127, questions 1 et 2.}\label{C04-M2} Distinguer l’ordonnée du point de contact et la pente de la droite tangente.
\begingroup\small
\mxnewpage\section*{QCM 1 : comprendre et choisir}\label{C04-Q1}
\mxCount{qcm1} questions pédagogiques à réponse unique. Choisir une réponse par question ; expliquer ensuite un distracteur. Ce format de chapitre ne reproduit pas le format officiel de l’EAM.
\begin{enumerate}
\item Taux de $x^2$ entre 1 et 2 :
\begin{enumerate}[label=\Alph*.,itemsep=2pt]
\item $2$
\item $3$
\item $1$
\item $4$
\end{enumerate}
\item Si $f(1)=4$, $f\prime(1)=2$, tangente :
\begin{enumerate}[label=\Alph*.,itemsep=2pt]
\item $y=2x+2$
\item $y=4x+2$
\item $y=2x+4$
\item $y=x+3$
\end{enumerate}
\item Le nombre dérivé est
\begin{enumerate}[label=\Alph*.,itemsep=2pt]
\item une fonction
\item une aire
\item une abscisse
\item un réel si f est dérivable en a
\end{enumerate}
\item Pour $x^2$ en 3, nombre dérivé :
\begin{enumerate}[label=\Alph*.,itemsep=2pt]
\item $9$
\item $3$
\item $6$
\item $12$
\end{enumerate}
\item L’approximation de $f(a+h)$ est
\begin{enumerate}[label=\Alph*.,itemsep=2pt]
\item $f(a)f\prime(a)$
\item $f(a)+hf\prime(a)$
\item $f(a)+f\prime(h)$
\item $hf(a)$
\end{enumerate}
\mxcounted{qcm1}
\end{enumerate}
\endgroup
\begingroup\small
\mxneed{18}\section*{QCM 2 : comprendre et choisir}\label{C04-Q2}
\mxCount{qcm2} questions pédagogiques à réponse unique. Choisir une réponse par question ; expliquer ensuite un distracteur. Ce format de chapitre ne reproduit pas le format officiel de l’EAM.
\begin{enumerate}
\item $\sqrt{x}$ en 0 est
\begin{enumerate}[label=\Alph*.,itemsep=2pt]
\item dérivable de dérivée 0
\item dérivable de dérivée 1
\item non dérivable
\item non définie
\end{enumerate}
\item $\dfrac{|h|}{h}$ pour $h<0$ vaut
\begin{enumerate}[label=\Alph*.,itemsep=2pt]
\item $-1$
\item $1$
\item $0$
\item $h$
\end{enumerate}
\item Une tangente de pente 0 est
\begin{enumerate}[label=\Alph*.,itemsep=2pt]
\item verticale
\item oblique
\item impossible
\item horizontale
\end{enumerate}
\item Position en m et temps en s : unité de la dérivée
\begin{enumerate}[label=\Alph*.,itemsep=2pt]
\item $\mathrm{m^2}$
\item $\mathrm{m\,s^{-1}}$
\item $\mathrm{s\,m^{-1}}$
\item $\mathrm{m}$
\end{enumerate}
\item Pour f dérivable en a, remplacer $f(a+h)$ par $f(a)+hf\prime(a)$ est en général
\begin{enumerate}[label=\Alph*.,itemsep=2pt]
\item une identité exacte
\item une définition de h
\item une approximation locale
\item toujours faux même pour une affine
\end{enumerate}
\mxcounted{qcm2}
\end{enumerate}
\subsection*{Vrai ou faux ? Justifier.}\label{C04-VF}
\begin{enumerate}
\item Une tangente ne peut jamais recouper une courbe.
\item La fonction valeur absolue est dérivable en 0.
\item Une fonction affine $mx+p$ a pour nombre dérivé m en tout réel.
\end{enumerate}
\begin{keypoint}[Bilan personnel]Je donne le domaine et h non nul ; je lis une pente ; j’utilise le point de tangence ; je distingue exact et approché.

Pour chaque objectif, indiquer : acquis sans aide ; réussi avec aide ; à retravailler. Écrire un exemple et une erreur à éviter.\end{keypoint}
\endgroup
\ifmxteacher\mxinclude{corriges/C04}\fi
